\documentclass[11pt]{article}
\usepackage[a4paper,margin=25mm]{geometry}
\usepackage{amsmath,amssymb,amsthm,mathtools}
\usepackage[T1]{fontenc}
\usepackage{lmodern,microtype}
\usepackage[hidelinks]{hyperref}
\newtheorem{theorem}{Theorem}
\newtheorem{proposition}[theorem]{Proposition}
\newtheorem{lemma}[theorem]{Lemma}
\newcommand{\C}{\mathbb C}
\newcommand{\R}{\mathbb R}
\newcommand{\A}{\mathcal A}
\newcommand{\J}{\mathcal J}
\DeclareMathOperator{\Spec}{Spec}
\DeclareMathOperator{\lc}{lc}
\title{Reflected relation norms and the exact common refinement of conormal algebras}
\author{Split-Zero research programme}
\date{13 September 2026}
\begin{document}
\maketitle

Fix an integer $k\ge1$, retain the original coordinate $S$ and the real
number $c=k/2$, and let $\chi\in\C[S]$ be monic of degree $q\ge1$.
The intended application uses the actual cyclic annihilator with every
root order retained. All the algebra below holds for every such $\chi$.
It completes the comparison raised by the full-jet transfer construction:
the squared norm of an original relation, its first conormal algebra,
and their common refinement have explicit maps even for a packet which
is not closed under reflection. The reflected factors and their phases
are displayed throughout.

\section{The original reflection and the norm polynomial}

Let a bar over a polynomial mean conjugation of its coefficients. Define
the conjugate-linear involution and its reflected monic polynomial by
\[
 f^\star(S)=\overline f(k-S),\qquad
 \chi^\dagger(S)=(-1)^q\chi^\star(S),\qquad
 N_\chi(S)=\chi(S)\chi^\dagger(S).
 \tag{RF.1}
\]
The factor $(-1)^q$ remains part of this definition. In particular
$\chi^\star=(-1)^q\chi^\dagger$ is the original involution image.
The change of polynomial coordinates and the full leading phase are
\[
 \Psi:\C[S]\xrightarrow{\sim}\C[u],\quad
 \Psi(f)(u)=f(c+iu),\qquad
 \Psi^{-1}(F)(S)=F((S-c)/i),\qquad
 \psi=i^{-q}\Psi(\chi).
 \tag{RF.2}
\]

\begin{proposition}
The polynomial $\chi^\dagger$ is monic of degree $q$,
$(\chi^\dagger)^\dagger=\chi$, and $N_\chi^\star=N_\chi$.
The original-coordinate identities are
\[
 \Psi(f^\star)=\overline{\Psi(f)},\qquad
 \Psi(\chi)=i^q\psi,\qquad
 \Psi(\chi^\dagger)=i^q\overline\psi,\qquad
 \Psi(N_\chi)=i^{2q}\psi\overline\psi.
 \tag{RF.3}
\]
Consequently the monic polynomial $\Pi=\psi\overline\psi$ has real
coefficients and satisfies $\Pi(u)=|\psi(u)|^2\ge0$ for real $u$.
The following three conditions are equivalent:
\[
 \chi^\dagger=\chi,\qquad
 \overline\chi(k-S)=(-1)^q\chi(S),\qquad
 \overline\psi=\psi.
 \tag{RF.4}
\]
Under these equivalent conditions, $N_\chi=\chi^2$ and $\Pi=\psi^2$.
\end{proposition}

\begin{proof}
The leading coefficient of $\overline\chi(k-S)$ is $(-1)^q$,
which proves monicity with its stated phase. Since $k$ is real,
$(f^\star)^\star=f$ and $(fg)^\star=f^\star g^\star$.
Applying these equations twice, including both real factors $(-1)^q$,
gives $(\chi^\dagger)^\dagger=\chi$ and
$N_\chi^\star=(-1)^{2q}\chi^\dagger\chi=N_\chi$.
Direct substitution gives
$\Psi(f^\star)(u)=\overline f(c-iu)=\overline{\Psi(f)}(u)$.
Conjugating $\Psi(\chi)=i^q\psi$ therefore gives
$\Psi(\chi^\star)=(-i)^q\overline\psi$.
Multiplication by $(-1)^q$ yields the third equation of (RF.3);
multiplication by $\Psi(\chi)$ yields the fourth. The two leading
coefficients in $\Pi$ equal one, and conjugating its product reverses
its factors, proving that its coefficients are real. Evaluation on
the real axis proves its value there. Finally (RF.3) and bijectivity
of $\Psi$ prove the equivalence (RF.4), including the factor $i^q$.
\end{proof}

The exact norm map can also be written without suppressing its phase.
For the original positive density $m_{h,k}(u)$ and any polynomial $Q$,
put $\widetilde Q=\Psi(Q)$. Polynomial integrability is supplied by the
original arithmetic source construction. Pointwise (RF.3) gives
\[
 \Psi(\chi Q)=i^q\psi\widetilde Q,
 \qquad
 \int_\R|\chi(c+iu)Q(c+iu)|^2m_{h,k}(u)\,du
 =\int_\R\Pi(u)|\widetilde Q(u)|^2m_{h,k}(u)\,du.
 \tag{RF.5}
\]
Thus multiplication by $\chi$ on the original polynomial source is
carried by $\Psi$ to $i^q$ times multiplication by $\psi$. Both
integrals use the original density and its literal mass. The equality
of integrands is $|i^q|^2|\psi\widetilde Q|^2
=\Pi|\widetilde Q|^2$, and proves the complete norm comparison.

\section{Full jets at the union of the original and reflected roots}

Let $a_\lambda$ be the order of $\lambda$ as a root of $\chi$, with
value zero at other points. The order of $\lambda$ in $\chi^\dagger$
is
\[
 b_\lambda=a_{k-\overline\lambda},\qquad
 r_\lambda=a_\lambda+b_\lambda.
 \tag{RF.6}
\]
To prove this, write $\chi(S)=\prod_\alpha(S-\alpha)^{a_\alpha}$
and substitute it in (RF.1). Each reflected factor is
$k-S-\overline\alpha=-(S-(k-\overline\alpha))$;
their $q$ minus signs cancel the displayed $(-1)^q$. Orders at
coincident roots add when the two polynomials are multiplied.
In particular $\sum_\lambda r_\lambda=2q$.

\begin{theorem}
There are explicit isomorphisms
\[
 \A_N:=\C[S]/(N_\chi)
 \xrightarrow{\ \Psi_N\ }\C[u]/(\Pi),\qquad
 [P]\longmapsto[P(c+iu)],
 \tag{RF.7}
\]
and, for any fixed ordering of the distinct roots in (RF.6),
\[
 \J_N:\A_N\xrightarrow{\sim}
 \bigoplus_{\lambda:r_\lambda>0}\C^{r_\lambda},\qquad
 [P]\longmapsto
       \bigl(P^{(d)}(\lambda)\bigr)_{\lambda,\,0\le d<r_\lambda}.
 \tag{RF.8}
\]
The product in each raw-derivative block is
$(v\cdot w)_d=\sum_{j=0}^d\binom dj v_jw_{d-j}$.
If $\zeta_\lambda=(\lambda-c)/i$, the same raw jets in the $u$
coordinate are exactly
\[
 (\Psi P)^{(d)}(\zeta_\lambda)=i^dP^{(d)}(\lambda).
 \tag{RF.9}
\]
Multiplication by $S$ in (RF.8) is the block map
$(M_Sv)_d=\lambda v_d+d v_{d-1}$, with $v_{-1}=0$.
\end{theorem}

\begin{proof}
Equation (RF.3) maps the ideal $(N_\chi)$ onto $(\Pi)$ because
$i^{2q}$ is an invertible constant. Applying $\Psi^{-1}$ proves
both well-definedness and inverse for (RF.7).
A polynomial has its first $r$ raw derivatives zero at $\lambda$
exactly when $(S-\lambda)^r$ divides it: its Taylor expansion is
$\sum_d P^{(d)}(\lambda)(S-\lambda)^d/d!$. For distinct roots
the corresponding factors are relatively prime. Repeated polynomial
division, or Bezout's identity applied successively to those factors,
shows that simultaneous divisibility is divisibility by their product
$N_\chi$. This proves the exact kernel in (RF.8).
Representatives of degree less than $2q$ are unique by monic division,
so the source and target of the induced injection both have dimension
$2q$; it is therefore onto. Leibniz's rule gives the stated block
product, including every binomial coefficient. The chain rule proves
(RF.9), and differentiating $SP$ gives
$(SP)^{(d)}(\lambda)=\lambda P^{(d)}(\lambda)
+dP^{(d-1)}(\lambda)$. These calculations also show that all maps
retain the entire nilpotent root order.
\end{proof}

There is a canonical original arithmetic quotient map
$\pi_N:\A_N\to\A_\chi:=\C[S]/(\chi)$.
Its exact kernel and its coefficient-module parametrization are
\[
 0\longrightarrow\C[S]/(\chi^\dagger)
 \xrightarrow{\ \iota_N\ }\A_N
 \xrightarrow{\ \pi_N\ }\A_\chi\longrightarrow0,
 \qquad\iota_N[f]=[\chi f].
 \tag{RF.10}
\]
Indeed $N_\chi\mid\chi f$ is equivalent to
$\chi^\dagger\mid f$ by cancellation in $\C[S]$.
The image consists precisely of the multiples of $\chi$ modulo
$N_\chi$, and reduction modulo $\chi$ is surjective. This proves
exactness, injectivity and all dimensions. The map $\iota_N$ is a
$\C[S]$-module map; the quotient maps are algebra homomorphisms.

\section{A common refinement with the original first conormal algebra}

Retain the original first conormal algebra
$\A_2=\C[S]/(\chi^2)$ and its exact sequence
\[
 0\longrightarrow\C[S]/(\chi)
 \xrightarrow{[f]\mapsto[\chi f]}\A_2
 \longrightarrow\A_\chi\longrightarrow0.
 \tag{RF.11}
\]
Let $g$ be the monic greatest common divisor of $\chi,\chi^\dagger$,
and put
\[
 d=\chi g,\qquad
 \ell=\chi^2\chi^\dagger/g.
 \tag{RF.12}
\]
These polynomials retain all original roots: at a point $\lambda$
their orders are respectively
$a_\lambda+\min(a_\lambda,b_\lambda)$ and
$a_\lambda+\max(a_\lambda,b_\lambda)$.

\begin{theorem}
The following sequence is exact as a sequence of $\C[S]$-modules:
\[
 0\longrightarrow\C[S]/(\ell)
 \xrightarrow{\ j\ }\A_2\oplus\A_N
 \xrightarrow{\ t\ }\C[S]/(d)\longrightarrow0,
 \quad
 j[P]=([P],[P]),\quad t([P],[Q])=[P-Q].
 \tag{RF.13}
\]
The first map is an algebra isomorphism onto the fibre product
\[
 \boxed{\C[S]/(\ell)\simeq
       \A_2\mathbin{\times}_{\C[S]/(\chi g)}\A_N.}
 \tag{RF.14}
\]
The projections of this fibre product followed by reduction modulo
$\chi$ agree. All the maps in (RF.10)--(RF.14) commute with the
original multiplication operator $M_S$.
\end{theorem}

\begin{proof}
In the polynomial ring, the intersection of two principal ideals
has at every root the maximum of their orders; their sum has the
minimum. The intersection assertion follows immediately from
divisibility at every root. For the sum, divide the two polynomials
by their monic greatest common divisor. The remaining factors are
coprime, and the Euclidean algorithm supplies a polynomial linear
combination equal to one. Multiplying back proves that their ideal
sum is the gcd ideal. Applying this to $\chi^2$ and $\chi\chi^\dagger$
gives the polynomials $\ell$ and $d$ in (RF.12).
Thus $j$ is well-defined and its kernel is zero. The difference
map $t$ is well-defined since $d$ divides both defining polynomials;
it is onto because $t([P],0)=[P]$.
If $P-Q$ belongs to their ideal sum, write
$P-Q=\chi^2 A+N_\chi B$. Then
$T=P-\chi^2 A=Q+N_\chi B$ maps to $([P],[Q])$ under $j$.
This proves $\ker t=\operatorname{im}j$ without choosing different
roots or changing any order. The image consists of exactly those
pairs with equal reduction modulo $d$, and its multiplication is
componentwise. Therefore the injection is the algebra isomorphism
(RF.14). Since $\chi$ divides $d$, the final reductions agree.
Multiplication by $S$ commutes with taking every polynomial residue,
with $P\mapsto\chi P$, and with the difference of representatives;
this proves the final assertion directly.
\end{proof}

On the two kernel modules over $\A_\chi$, the comparison map has
an especially explicit description. Reduction to $\C[S]/(\chi g)$
sends $[\chi f]$ from either algebra to the same residue class.
Its image is $(\chi)/(\chi g)$, canonically isomorphic to
$\C[S]/(g)$ by $[f]\mapsto[\chi f]$. Thus the exact diagram of
kernel maps is the pair of quotient maps
\[
 \C[S]/(\chi)\longrightarrow\C[S]/(g)
 \longleftarrow\C[S]/(\chi^\dagger),
 \tag{RF.15}
\]
with their literal multiplication-by-$\chi$ parametrizations in
(RF.11) and (RF.10). This is the relation between the first
conormal kernel and the kernel recorded by the relation norm.

\section{Square-zero structure, reflected sectors, and supported zero}

The kernel in (RF.11) has square zero, because products of its
representatives contain $\chi^2$. The kernel in (RF.10) has square
zero exactly when $\chi^\dagger=\chi$. To prove necessity, its
element $[\chi]$ would have square zero, requiring
$\chi\chi^\dagger\mid\chi^2$, and hence
$\chi^\dagger\mid\chi$. Equal degree and monicity force equality.
Sufficiency follows from the same divisibility. The exact common
refinement for both cases remains (RF.13)--(RF.15).

For an explicit complementary case suppose $g=1$. Choose the
polynomials $u,v$ supplied by the Euclidean algorithm with
$u\chi+v\chi^\dagger=1$. Then the algebra map
\[
 \A_N\xrightarrow{\sim}\A_\chi\times\A_{\chi^\dagger},
 \quad[P]\longmapsto([P]_{\chi},[P]_{\chi^\dagger}),
 \qquad e_{\rm ref}=[u\chi]_{N_\chi}
 \tag{RF.16}
\]
is an isomorphism. Its kernel before factoring is the intersection
$(\chi)\cap(\chi^\dagger)=(N_\chi)$, and for arbitrary residues
represented by $a,b$, the polynomial
$a v\chi^\dagger+b u\chi$ has those two residues. This proves
surjectivity and the inverse formula. The element $e_{\rm ref}$
has coordinates $(0,1)$, proving $e_{\rm ref}^2=e_{\rm ref}$.
Its ideal is the kernel of $\pi_N$ and is isomorphic as a unital
algebra with its local unit $e_{\rm ref}$ to
$\A_{\chi^\dagger}$. In these coordinates $M_S$ acts by the
original residue of $S$ in each factor. Equations (RF.14)--(RF.16)
therefore calculate the reflected component and its exact map to
the original arithmetic quotient.

To state the Split-Zero observation explicitly, attach any original
support label $\lambda_0$ to each of these vector spaces. A linear
map $L:E\to F$ has the support-preserving lift
\[
 \widetilde L(\tau)=\tau,\qquad
 \widetilde L(\lambda_0,x)=(\lambda_0,Lx).
 \tag{RF.17}
\]
Composition satisfies $\widetilde{L_2L_1}
=\widetilde L_2\widetilde L_1$ by evaluation on $\tau$ and on
each labelled vector. Addition and scalar action within a fibre
are preserved because $L$ is linear. Every vector of the kernels
in (RF.10), (RF.11), or (RF.13) maps to the receiving supported
zero $(\lambda_0,0)$; the input and all its pre-quotient coordinates
remain defined. In particular, the nonzero algebra idempotent
$e_{\rm ref}$ in (RF.16) has zero amplitude in $\A_\chi$, and
its labelled image is $(\lambda_0,0)$. These two values and the
external point $\tau$ have been related by the explicit quotient,
rather than identified in the source.

The first conormal algebra used here is exactly $\C[S]/(\chi^2)$.
If a tensor construction supplies a different annihilator, its
divisibility maps require that actual polynomial. No replacement
of a tensor annihilator is made by (RF.12): (RF.13) is the complete
comparison for the two polynomials explicitly stated here.

\section{Provenance and verification scope}

This derivation follows the explicit dagger-stability correction
to the full-jet transfer note at revision
\texttt{dfcbba5cbf7fec8c9301fe242c13e741d762013e}.
The endpoint product continuation retains the full transfer
calculation for $\Pi=\psi\overline\psi$. The present note proves
the original-coordinate helper algebra, both kernel modules,
and their exact gcd/lcm refinement, including their raw jets and
original action. These are polynomial algebra proofs with all
domains and maps displayed. The source norm identity (RF.5)
uses the unchanged arithmetic density; it does not estimate the
growth of any quotient volume or establish a zero-location claim.
\end{document}
